\section{Introduction}
An experiment is used through a causal observer. Information acquired
at one time may be lost at a later checkpoint, and a decision rule
that is economical after training may be expensive to train. The
statistical value of an observer therefore depends on its entire
register profile, not only on the dimension of its last statistic.
A further difficulty is that two feasible observers need not be
mixable at the same profile: the variable that selects an observer
has to be retained wherever their future actions differ.

We study this problem from positive experimental kernels, preparations,
causal interventions and bounded executable tests. These objects do
not presuppose a convex parameter image, an information metric, or a
finite-dimensional closure of the underlying dynamics. The present
paper develops their finite resource layer and its comparison maps.
The word ``finite'' refers to a declared protocol and its observable
alphabets, not to a removal of the microscopic or limiting questions
in the larger theory.

\subsection{The full-profile problem}
Fix a finite schedule $\sigma$, a register profile $\mathbf K$, a
family of environments $B$, and a terminal payoff $G\in[-1,1]$. Let
\[
 V^{\rm clk}_{\sigma,\mathbf K}(B,G)
       =\sup_A\inf_{b\in B}\E_b^A G,
\]
where every persistent data-dependent variable of $A$ passes through
the specified registers. The clocked model allows a predetermined
transition table at each epoch; its clock carries no discarded data.
An autonomous controller instead reuses its declared transition rows
and must retain any phase variable needed to select a different row.
Refining an atomic operation into input-bit or random-bit epochs
changes $\sigma$ and its resource problem. These conventions are part
of the theorem statements.

The central object is a compatible occupation array. Its coordinates
record the probability mass entering and leaving each observable
controller row, under every environment and hidden history. Ordinary
flow conservation is necessary but not sufficient. All hidden histories
that present the same observable row must use the same stochastic
transition. This requirement is a collection of quadratic proportionality
equations. Unlike a division-based parametrization, these equations
also retain every zero-occupation boundary case.

\begin{theorem}[Compatible statistical resource frontiers]
\label{thm:v23-main}
The following statements hold for the indicated finite protocols.

\smallskip\noindent
\textup{(i)} For a finite environment family, compatible occupation
arrays are exactly the arrays of all stochastic machines of the full
profile. Their minimax value is attained. It equals the decreasing
limit of polynomial sum-of-squares upper certificates, and every
feasible array reconstructs a machine. Shared-row equations give the
corresponding autonomous characterization. Finite-precision controller
witnesses and compact-family kernel nets give matching two-sided
approximations.

\smallskip\noindent
\textup{(ii)} For binary, conditionally independent finite observations,
possibly with a finite controlled choice of experiment, the optimal
clocked Bayesian value at any finite register profile is given by a
recursion of compatible interval partitions of incoming posterior
atoms. For uncontrolled observations with one- or two-state registers,
a receiver-curve recursion gives the exact clocked minimax value at
every finite horizon. These are optimizations over all stochastic
machines, not implementation bounds for a chosen statistic.

\smallskip\noindent
\textup{(iii)} With three independent binary reports of error probability
$p\in(0,1/2)$, put $d=p(1-p)(1-2p)$. At the atomic profile $(2,2,2)$
the exact fair-prior and minimax successes are respectively
$1-p+d/2$ and $1-p+d/3$. At profile $(2,3,2)$ both equal $1-p+d$.
A free mixture of the two Bayesian optimal machines would raise the
first profile's minimax value by $d/6$, so that convexification is
strictly incorrect for that resource problem.

\smallskip\noindent
\textup{(iv)} A fixed $K$-response calibration can be acquired from
$M=2^s$ target preparations, with uniform error at most $\rho$ except
with probability $\beta$, whenever
$M\ge(2\rho^2)^{-1}\log(2K/\beta)$. Its $(M+1)^K$-state counter
vector is retained and priced through the subsequent auditor. The
resulting response-cover constructions and their physical consumers
belong to the same full-profile classes as the occupation converse.
\end{theorem}
\begin{proof}
Part (i) is Theorems~\ref{thm:v23-occupation} and
\ref{thm:v23-dual} and Corollary~\ref{cor:v23-compact}. Part (ii) is
Theorems~\ref{thm:v23-interval} and \ref{thm:v23-minimax-recursion}.
Part (iii) is Theorem~\ref{thm:v23-bsc} and
Corollary~\ref{cor:v23-mixing-gap}. Part (iv) is
Theorem~\ref{thm:v23-calibration} and
Corollary~\ref{cor:v23-reset-consumer}.
\end{proof}

The certificate hierarchy in (i) uses a classical positivity theorem.
Its role is to give an upper bound on every complete machine, paired
with the reverse reconstruction from feasible occupations. It is not
an efficiency or finite-termination assertion. The noisy recursions in
(ii) and the closed formulas in (iii) supply structural information
that a generic polynomial encoding alone would not provide. In
particular, the distinction between Bayesian and minimax objectives
survives in an experiment whose two laws have common support at every
observation.

\subsection{Comparison with existing methods}
Comparison of experiments is classical
\cite{Blackwell1953,LeCam1964,Torgersen1991}. Filtered randomization is
also an existing subject. Weisshaupt's filtered criterion
\cite[Lemma 1 and Theorem 5]{Weisshaupt2002} concerns compact convex
classes of positive operators and a risk criterion for filtered
randomization. Our feasible object is a specified finite register
execution, whose behavior set need not be convex. We also cite
Norberg~\cite{Norberg2002}; the original proof text has not been obtained
for a definitive theorem-level comparison, and no priority claim rests
on that gap in access.

Finite-memory statistical learning and the role of time-dependent
transition rules are treated by Cover and by Hellman and Cover
\cite{Cover1969,HellmanCover1970}. Likelihood-ratio quantizer optimality
is a classical mechanism~\cite[Section III]{Tsitsiklis1993}. We use it
both for compatible posterior partitions and for binary-channel
domination. The receiver-curve and equal-error arguments are not
introduced as new statistical principles. Their role here is to solve
one full-profile optimization, including its stochastic minimax part
and the selector cost exposed by the explicit three-report frontier.

Polynomial descriptions of partially observed controller occupations
are likewise established. M\"uller and Mont\'ufar
\cite[Example 15, Theorem 16 and Remark 19]{MullerMontufar2022} give
semialgebraic feasible-frequency descriptions and polynomial planning
for infinite-horizon memoryless POMDPs, with their stated positivity
assumptions. The finite-protocol formulation here makes the observable
row equivalences, changing register alphabets and zero-mass
reconstruction explicit. Its upper-certificate completeness applies
Putinar's theorem and its rational version
\cite{Putinar1993,Powers2009}. The combined contribution is the exact
compatible frontier, its noisy evaluations, and its connections to
fully priced experimental consumers; neither polynomial optimization
nor the vanishing-minor identity is claimed as a new invention.

\subsection{Organization and scope of the applications}
Sections~\ref{sec:v23-occupations}--\ref{sec:v23-consumers} contain the
principal proof chain. The marked continuation and decision spectra
then identify the underlying experimental interfaces. The finite
precision, response-cover tournaments, multicut revelation formula,
weighted residuals, nonlinear tests and adaptive tests are developed
with their proofs. The microscopic moving-collision construction and
the bounded Gaussian comparison produce physical consumers of the
same interfaces. The original short, polynomial-horizon and long-clock
physical resource points are all retained.

Those three points are constructions, not an exact physical Pareto
curve. The general converse now applies to each of their fixed
protocols, but its numerical evaluation across all physical budgets
is a separate task. Likewise a confidence lower bound at a specified
retained-coin interface is not automatically a lower bound for an
auditor that uses richer raw observations. These distinctions are
mathematical typing conditions on the applications, not changes to
the physical or operator targets.

The article retains the full development of its supporting results;
versioned derivation and review records are kept outside the exposition.
The dependency appendix records the namespaced eleven-component
program. The primary A2 geometric chain remains independent, while
B4 kinetic and broader C2 operator conclusions require their own stated
hypotheses and proofs. No artificial dependence on a finite-state
theorem is imposed on those branches of the program.
